\section{Retaining more variance than the degree cost}
\label{sec:observation}

We now construct the observation rule needed for amplification.  Revealing
any fixed set of $m$ of $m+1$ independent centered variance-one inputs retains variance
$m$ from their sum: the unobserved input accounts for the remaining unit.
Our rule will retain strictly more than $m$, while every output indicator
is a linear combination of functions of at most $m$ inputs.  These are
the two conclusions required by Proposition~\ref{prop:amplification},
with $r=m+1$ and $d=m$.

The rule usually leaves one input unobserved.  On one specially chosen
event, it instead reports only that the event occurred.  This improves
the prediction error if the sum of all the inputs is sufficiently close
to one fixed number on that event.  We first describe the rule and
calculate its error for arbitrary independent inputs.  We will then
choose its parameters using Gaussian inputs.

\subsection{A reporting rule and its prediction error}

Let $m\ge3$, let $I_1,\ldots,I_m$ be intervals, and let
$\iota:\mathbb R\to\{1,\ldots,m\}$ be a selector with finitely many
cut points.  All endpoint and tie conventions are part of the rule.
Write its inputs as $(u,z_1,\ldots,z_m)$ and their sum as
$L=u+\sum_{j=1}^m z_j$.  The coordinate $u$ selects the leaf
$z_{\iota(u)}$.  Define
\[
 B=\{z_j\in I_j\text{ for every }j\},\qquad
 P=\{z_j\in I_j\text{ for every }j\ne\iota(u)\},\qquad
 A=P\setminus B.
\]
Thus $A$ is the event that the selected leaf alone fails its interval
test.  Since $B\subseteq P$, the sets $A$, $B$, and $P^c$ partition
the input space.  Define the tagged observation $\mathcal W$ by
\[
\begin{array}{c|l}
 \text{Event}&\text{Information reported}\\ \hline
 A&\text{the tag }A\text{ only}\\
 B&(B,z_1,\ldots,z_m)\\
 P^c&(P^c,u,\iota(u),(z_j)_{j\ne\iota(u)}).
\end{array}
\]
On $B$, the central input is omitted.  On $P^c$, the selected leaf
is omitted.  Although the exceptional report on $A$ can depend on all
$m+1$ inputs, its indicator has the pointwise identity
\begin{equation}\label{eq:indicator}
 \mathbf1_A=
 \sum_{i=1}^m\mathbf1_{\{\iota(u)=i\}}
                  \prod_{j\ne i}\mathbf1_{I_j}(z_j)
 -\prod_{j=1}^m\mathbf1_{I_j}(z_j).
\end{equation}
Each summand uses only $m$ coordinates.  This cancellation is what
allows the exceptional report to have the same degree cost as the
ordinary reports.

\begin{lemma}\label{lem:report}
For every product of finite input alphabets, the rule $\mathcal W$
has finite range, and each output indicator is a finite linear
combination of functions of at most $m$ input coordinates.

For independent centered variance-one inputs, whose marginal laws need
not agree, and every fixed $b\in\mathbb R$, put
\[
 \Delta_b=\mathbb E[\mathbf1_A(1-(L-b)^2)].
\]
Then
\begin{equation}\label{eq:retained}
 \operatorname{Var}(\mathbb E[L\mid\mathcal W])\ge m+\Delta_b.
\end{equation}
\end{lemma}

\begin{proof}
Identity~\eqref{eq:indicator} gives the required decomposition
for the $A$ report.  An attained $B$ report fixes the $m$ leaf values
and imposes no condition on $u$.  An attained $P^c$ report fixes $u$,
its selected index $i$, and the $m-1$ unselected leaves.  At least one
of these reported leaves fails its test; hence the reported values
already force $P^c$, whatever the omitted value $z_i$ may be.  Its
indicator therefore depends on only the $m$ reported coordinates.
Unused output labels have zero indicators.  This proves the first
assertion on every input, including interval endpoints.

To estimate the retained variance, predict $L$ by $b$ on $A$, by
$\sum_jz_j$ on $B$, and by $u+\sum_{j\ne\iota(u)}z_j$ on $P^c$.
This predictor is a function of the reported information.  On $B$
its error is $u$; independence gives
$\mathbb E[\mathbf1_Bu^2]=\mathbb P(B)$.
For the last type of report, fix an index $i$ and let
$E_i=P^c\cap\{\iota(u)=i\}$.  This event and its reported data
depend only on $u$ and the leaves other than $z_i$.  Consequently
\[
 \mathbb E[\mathbf1_{E_i}z_i]=0,
 \qquad
 \mathbb E[\mathbf1_{E_i}z_i^2]=\mathbb P(E_i).
\]
Summing the squared errors over the three types of report gives
\[
 \mathbb P(A^c)+\mathbb E[\mathbf1_A(L-b)^2]=1-\Delta_b.
\]
Conditional expectation has no larger mean squared error
\cite[Theorems~4.1.13 and 4.1.15]{Durrett}.  Since $L$ is centered and
$\operatorname{Var}(L)=m+1$, the conditional-variance identity gives
\eqref{eq:retained}.
\end{proof}

The remaining task is now precise: choose the intervals, selector, and
constant $b$ so that $\Delta_b>0$ for Gaussian inputs.  A strict
positive margin will survive the normal approximation used in
amplification.  The next calculation produces that margin by subtracting
the contribution of $B$ from that of $P$ in
\eqref{eq:indicator}.

\subsection{Choosing the intervals in the Gaussian model}

We will obtain a positive predictor deficit and also record the
conditional variance on $A$.  The latter describes the event's gain
without specifying a predictor: on $A$, the sum varies less than one
unobserved variance-one input.

\begin{lemma}\label{lem:gaussian}
There are a finite odd integer $m\ge3$, closed intervals
$I_1,\ldots,I_m$, and a selector $\iota$ with finitely many cut points
such that the following holds for independent standard normal inputs.
Let $\beta_j$ be the probability of $I_j$, and let $\mu_j,s_j^2$ be
the conditional mean and variance on $I_j$.  Define
\[
 p_* =\prod_{j=1}^m\beta_j,\qquad
 S=\sum_{j=1}^m s_j^2,\qquad
 b=\sum_{j=1}^m\mu_j-1.
\]
Then $p_*>0$, $S<1/4$, and
\begin{equation}\label{eq:deficit}
 d_A:=\mathbb E[\mathbf1_A(1-(L-b)^2)]>\frac34p_*.
\end{equation}
The boundary of $A$ lies in a finite union of coordinate hyperplanes.
In particular, $\mathbb P(A)>0$ and $\operatorname{Var}(L\mid A)<1$.
\end{lemma}

\begin{proof}
The shift by one in $b$ determines where to place the intervals.
Conditional on $u=t$ and $P$, the mean of $L-b$ is
$1+t-\mu_{\iota(t)}$, so selecting an interval near $1+t$ makes
this mean small.  Conditional on $B$ alone, with $u$ still random,
its mean is instead one.  Subtracting the contribution of $B$ will
turn that fixed squared mean into the positive term in $d_A$.

We let odd $m$ tend to infinity to choose the parameters, and then fix
one finite value.  Set
\[
 R=\sqrt{8\log m},\qquad h=m^{-3/2},\qquad
 a_i=1-R+\frac{2R(i-1)}{m-1},\qquad I_i=[a_i-h,a_i+h].
\]
For $|t|\le R$, select a center $a_{\iota(t)}$ nearest to $1+t$,
resolving ties by the smaller index.  For $|t|>R$, use the middle
index, whose center is $1$.  These choices define the rule on the
whole real line.  Its finitely many interval endpoints and selector
cut points also give the asserted boundary property.

Write $\phi(t)=(2\pi)^{-1/2}e^{-t^2/2}$.  Every $\beta_i$ lies in
$(0,1)$.  Condition first on $u=t$ and put $i=\iota(t)$.  The event
$P$ then has probability $p_*/\beta_i$; it restricts the leaves
$j\ne i$ to their intervals and leaves $z_i$ unrestricted.  Under
this conditioning, $L-b$ has mean $1+t-\mu_i$ and variance
$1+\sum_{j\ne i}s_j^2$.  Thus
\[
 \mathbb E[\mathbf1_P(1-(L-b)^2)]=-p_*J,
 \qquad
 J=\int_{\mathbb R}
 \left((1+t-\mu_{\iota(t)})^2+
                   \sum_{j\ne\iota(t)}s_j^2\right)
                 \frac{\phi(t)}{\beta_{\iota(t)}}\,dt.
\]
On $B$, with $u$ again random, $L-b$ has mean $1$ and variance
$1+S$.  Therefore
\[
 \mathbb E[\mathbf1_B(1-(L-b)^2)]=-p_*(1+S).
\]
Since $\mathbf1_A=\mathbf1_P-\mathbf1_B$, we obtain the exact identity
\begin{equation}\label{eq:gaussian-identity}
 d_A=p_*(1+S-J).
\end{equation}
The positive constant $1$ comes from the mean of $L-b$ on $B$.
It remains to make the nonnegative error $J$ small.  Nearest-center
selection controls its central part; the fallback to the middle
interval controls the tails.

The elementary interval estimates are
\[
 |\mu_i-a_i|\le h,\qquad s_i^2\le h^2,\qquad
 \beta_i\ge 2h(2\pi)^{-1/2}
              \exp\!\left(-\frac{(|a_i|+h)^2}{2}\right).
\]
For $|t|\le R$, the grid spacing and the selector give
\[
 |1+t-\mu_{\iota(t)}|\le\frac{R}{m-1}+h,
 \qquad |a_{\iota(t)}|+h\le |t|+2
\]
when $m$ is sufficiently large.  Hence
$\phi(t)/\beta_{\iota(t)}\le(e^2/2h)e^{2R}$ on this range, and
\begin{align*}
 J_{\mathrm{central}}
 &\le\frac{e^2R}{h}e^{2R}
       \left(\left(\frac{R}{m-1}+h\right)^2+mh^2\right)\\
 &=O\!\left(R(R^2+1)e^{2R}m^{-1/2}\right)=o(1).
\end{align*}
Here $e^{2R}=m^{o(1)}$, while $R$ is a power of $\log m$.

On $|t|>R$, the selected interval is centered at $1$.  For $h\le1$
its probability is at least $2h\phi(2)$, and the numerator in $J$
is $O(1+t^2)$.  Integration by parts gives
\[
 \int_{|t|>R}(1+t^2)\phi(t)\,dt
       =O((R+1)e^{-R^2/2}).
\]
Consequently
\[
 J_{\mathrm{tail}}
 =O(h^{-1}(R+1)e^{-R^2/2})
 =O((R+1)m^{-5/2})=o(1).
\]
Also $S\le mh^2=m^{-2}\to0$.  Fix a sufficiently large finite odd
$m$ so that $J<1/4$ and $S<1/4$.  Identity~\eqref{eq:gaussian-identity}
now proves \eqref{eq:deficit}, with $p_*>0$.

Finally, write $p_A=\mathbb P(A)$.  Since $d_A>0$, we have $p_A>0$,
and
\[
 d_A=p_A\bigl(1-\operatorname{Var}(L\mid A)
                   -(\mathbb E[L\mid A]-b)^2\bigr)>0.
\]
This also proves $\operatorname{Var}(L\mid A)<1$.
\end{proof}

\subsection{Transfer to finite inputs and amplification}

Fix all the Gaussian choices from Lemma~\ref{lem:gaussian},
including $m$, the intervals, the selector, $b$, and $p_*>0$.
Set $\eta=p_*/2$.  If centered variance-one laws $Z_n$ converge
weakly to a standard normal and have uniformly bounded fourth moments,
apply Lemma~\ref{lem:transfer} to the fixed function
\[
 g(u,z_1,\ldots,z_m)=\mathbf1_A(1-(L-b)^2).
\]
It has quadratic growth and discontinuities only on the finitely many
hyperplanes from Lemma~\ref{lem:gaussian}; these have Gaussian
measure zero.  For independent copies of $Z_n$, therefore,
$\mathbb Eg\to d_A>3p_*/4$, and eventually $\mathbb Eg>\eta$.
Lemma~\ref{lem:report} then gives
\[
 \operatorname{Var}\!\left(
   \mathbb E\!\left[\sum_{i=1}^{m+1} Z_{n,i}
                   \,\middle|\,\mathcal W(\mathbf Z_n)\right]\right)
 >m+\eta.
\]
Together with the pointwise output-indicator decomposition in that
lemma, this verifies both hypotheses of
Proposition~\ref{prop:amplification} with
$(r,d,\eta)=(m+1,m,p_*/2)$.

The Gaussian parameters and $\eta$ remain fixed throughout the iteration.
Proposition~\ref{prop:amplification} therefore increases the ratio of
retained variance to cell degree bound by at least $1+\eta/m>1$ at
each stage, using a sufficiently large finite copy count chosen for the
current observation. Starting from one revealed sign, finitely many
stages make this ratio exceed $4C^2$ for any prescribed $C>0$.
Lemma~\ref{lem:readout} then supplies a nonconstant Boolean function
with signed singleton sum greater than $C\sqrt{\deg(f)}$, proving
Theorem~\ref{thm:main}.
